\subsection{HAAR MEASURE ON A LIE GROUP}

Let G be a Lie group of finite dimension n.~Then \(\bigwedge^{n}(\mathrm{T}_{e}(\mathrm{G}))\) is of dimension 1. Hence (no. 13) the vector space S of left invariant differential forms of degree n on G is of dimension 1. Let \((\omega_{1},\ldots,\omega_{n})\) be a basis of the space of left invariant differential forms of degree 1 on G; then \(\omega_{1}\wedge\omega_{2}\wedge\cdots\wedge\omega_{n}\) is a basis of S.

PROPOSITION 54. Let G be a Lie group of finite dimension n, \(\omega\) a left invariant differential form of degree n on G and \(\phi\) an endomorphism of G. Then

\[
\phi^ {*} (\omega) = (\det L (\phi)) \omega .
\]

We write \(\mathrm{L}(\phi) = u, \omega_e = f\) and \(\phi^*(\omega)_e = g\) . For all \(x_1, \ldots, x_n\) in \(\mathrm{L}(\mathrm{G})\) ,

\[
g (x _ {1}, \dots , x _ {n}) = f (u x _ {1}, \dots , u x _ {n}) = (\det u) f (x _ {1}, \dots , x _ {n})
\]

and hence \(\phi^{*}(\omega)_{e} = \det L(\phi). \omega_{e}\) . On the other hand, if \(g \in G\) ,

\[
\phi \circ \gamma (g) = \gamma (\phi (g)) \circ \phi
\]

and hence \(\gamma(g)^{*}\phi^{*}(\omega)=\phi^{*}(\omega)\) . Thus \(\phi^{*}(\omega)\) is left invariant, whence the proposition.

COROLLARY. For all \(g \in G\) ,

\[
\delta (g) ^ {*} \omega = (\det \operatorname{Ad} g) \omega .
\]

\[
\delta (g) ^ {*} \omega = \delta (g) ^ {*} \gamma (g) ^ {*} \omega = (\operatorname{Int} g) ^ {*} \omega \text {   and   } L (\operatorname{Int} g) = \operatorname{Ad} g.
\]

Let G be a locally compact group and \(\phi\) an endomorphism of G. Suppose that there exist open neighbourhoods V, \(V'\) of e such that \(\phi(V) = V'\) and \(\phi \mid V\) is a local isomorphism of G into G. Let \(\mu\) be a left Haar measure on G. By Integration, Chapter VII, §1, Corollary to Proposition 9, there exists a unique number a \textgreater{} 0 such that \(\phi(\mu \mid V) = a^{-1}\mu \mid V'\) . Clearly a is independent of the choice of V, \(V'\) and \(\mu\) . It is called the modulus of \(\phi\) and is denoted by mod \(_{G}\) \(\phi\) or simply mod \(\phi\) . When \(\phi\) is an automorphism of G, we recover Definition 4 of Integration, Chapter VII, §1.

PROPOSITION 55. Suppose that K is locally compact. Let \(\mu\) be a Haar measure on the additive group of K. Let G be a Lie group of finite dimension \(n\) .

\begin{enumerate}
\def\labelenumi{(\roman{enumi})}
\item
  Let \(\omega\) be a non-zero left invariant differential form of degree \(n\) on G. Then the measure mod \((\omega)_{\mu}\) (Differentiable and Analytic Manifolds, R, 10.1.6) is a left Haar measure on G. If K = R and G has the orientation defined by \(\omega\) , the measure defined by \(\omega\) (Differentiable and Analytic Manifolds, R, 10.4.3) is a left Haar measure on G.
\item
  Let \(\phi\) be an étale endomorphism of G. Then mod \(\phi = \text{mod}\det L(\phi)\) .
\item
  is obvious. Let \(\mathrm{V}, \mathrm{V}'\) be open neighbourhoods of \(e\) such that \(\phi(\mathrm{V}) = \mathrm{V}'\) and \(\phi \mid \mathrm{V}\) is a local isomorphism of \(\mathrm{G}\) into \(\mathrm{G}\) . Then
\end{enumerate}

\[
\begin{array}{r l r l} \phi^ {- 1} (\mathrm{mod} (\omega) _ {\mu}   |   \mathrm{V} ^ {\prime}) & = \mathrm{mod} (\phi^ {*} (\omega)) _ {\mu}   |   \mathrm{V}) & & \text {by transport of structure} \\ & = \mathrm{mod} (\det \mathrm{L} (\phi) \omega   |   \mathrm{V}) _ {\mu} & & \text {(Proposition 54)} \\ & = \mathrm{mod} \det \mathrm{L} (\phi) (\mathrm{mod} (\omega) _ {\mu}   |   \mathrm{V}) \end{array}
\]

whence mod \(\phi = \mathrm{mod}\det \mathrm{L}(\phi)\) by definition of mod \(\phi\) .

COROLLARY. For all \(g \in \mathbf{G}\) , \(\Delta_{\mathbf{G}}(g) = (\text{mod det } \operatorname{Ad} g)^{-1}\) . In particular, for \(\mathbf{G}\) to be unimodular, it is necessary and sufficient that mod det \(\operatorname{Ad} g = 1\) for all \(g \in \mathbf{G}\) .

\[
\begin{array}{r l} \Delta_ {\mathbf {G}} (g) & = (\text {mod} \operatorname{Int} g) ^ {- 1} \quad \text {(Integration, Chapter VII, § 1, formula (33))} \\ & = (\text {mod} \det \operatorname{L} (\operatorname{Int} g)) ^ {- 1} \quad \text {(Proposition 55)} \\ & = (\text {mod} \det \operatorname{Ad} g) ^ {- 1}. \end{array}
\]

Remark. Preserving the hypotheses and notation of Proposition 52, suppose that K is locally compact. Let \(\mu\) be the measure

\[
\operatorname{mod} \det (\mathrm{D} _ {n + i} m _ {k} (x _ {1}, \dots , x _ {n}, x _ {1}, \dots , x _ {n})) _ {1 \leqslant i, k \leqslant n} d x _ {1} \dots d x _ {n}
\]

on \(\psi (\mathbf{U})\) . Then \(\psi^{-1}(\mu)\) is the restriction to U of a Haar measure on G.

PROPOSITION 56. Let G be a Lie group of finite dimension n, H a p-dimensional Lie subgroup and X the Lie homogeneous space G/H. Suppose that

\[
\det \operatorname{Ad} _ {\mathrm{L} (\mathrm{G})} h = \det \operatorname{Ad} _ {\mathrm{L} (\mathrm{H})} h
\]

for all \(h\in \mathbf{H}\) . Then:

\begin{enumerate}
\def\labelenumi{(\roman{enumi})}
\item
  The differential forms of degree \(n - p\) on \(\mathbf{X}\) which are invariant under \(\mathbf{G}\) are analytic.
\item
  The vector space of these forms is of dimension 1.
\item
  If \(\omega\) is such a non-zero form and K is locally compact, mod \((\omega)_{\mu}\) is a non-zero measure on X which is invariant under G.
\end{enumerate}

By § 1, no. 8, Examples, \(\mathrm{Alt}^{n - p}(\mathrm{TX},\mathrm{K})\) is an analytic vector G-bundle. Let \(x_0\) be the canonical image of \(e\) in \(\mathbf{X}\) ; its stabilizer is H. The fibre of \(\mathrm{Alt}^{n - p}(\mathrm{TX},\mathrm{K})\) at \(x_0\) is \(\bigwedge^{n - p}\mathrm{T}_{x_0}(\mathbf{X})^*\) and \(\mathrm{T}_{x_0}(\mathbf{X})\) is canonically identified with \(\mathrm{L}(G) / L(\mathrm{H})\) . If \(h\in \mathrm{H}\) , the automorphism \(\tau_h\) of \(\mathbf{X}\) defined by \(h\) is derived when passing to the quotient from the automorphism \(g\mapsto hgh^{-1}\) of G. Hence the automorphism \(\mathrm{T}_{x_0}(\tau_h)\) is derived when passing to the quotient from \(\mathrm{Ad}_{\mathrm{L(G)}}(h)\) . As

\[
\det \mathrm{Ad} _ {\mathrm{L(G)}} h = (\det \mathrm{Ad} _ {\mathrm{L(H)}} h). (\det \mathrm{T} _ {x _ {0}} (\tau_ {h})),
\]

the hypothesis implies that \(\det \mathrm{T}_{x_0}(\tau_h) = 1\) . Thus, every element of \(\bigwedge^{n - p}\mathrm{T}_{x_0}(\mathbf{X})^*\) is invariant under H. Then (i) and (ii) follow from § 1, no. 8, Corollary 1 to Proposition 17 and (iii) is obvious.

The existence of a non-zero positive measure on X invariant under G follows from Integration, Chapter VII, § 2, Corollary 2 to Theorem 3, for the hypothesis of Proposition 56 implies \(\Delta_{\mathbf{G}}|H = \Delta_{\mathbf{H}}\) (Corollary to Proposition 55).

PROPOSITION 57. Let G be a Lie group of finite dimension n.~Choose a basis for \(\bigwedge^{n}\mathrm{T}_{e}(\mathrm{G})^{*}\) ; by means of the right (resp. left) trivialization of \(\bigwedge^{n}\mathrm{T}(\mathrm{G})^{*}\) , we can identify this vector bundle with the trivial vector bundle \(\mathbf{G} \times \mathbf{K}\) , so that the transpose of a scalar differential operator is identified with a scalar differential operator.

Then, if \(u \in \mathrm{U}(\mathrm{G})\) , the transpose of \(\mathrm{L}_u\) (resp. \(\mathbf{R}_u\) ) is \(\mathrm{L}_u^\vee\) (resp. \(\mathbf{R}_u^\vee\) ).

We shall consider the case where \(\bigwedge^{n}\mathrm{T}(\mathrm{G})^{*}\) has been trivialized using a right invariant form \(\omega\) .

Suppose that the proposition has been proved for elements \(u_{1}, u_{2}\) of U(G). Then,

\[
\begin{array}{l l} ^ {t} (\mathrm{L} _ {u _ {1} * u _ {2}}) = ^ {t} (\mathrm{L} _ {u _ {1}} \circ \mathrm{L} _ {u _ {2}}) & \text {(Proposition 23)} \\ = ^ {t} (\mathrm{L} _ {u _ {2}}) \circ^ {t} (\mathrm{L} _ {u _ {1}}) & \text {(Diff. \& Anal. Man., R, 14.3.3)} \\ = \mathrm{L} _ {u _ {2}} ^ {\vee} \circ \mathrm{L} _ {u _ {1}} ^ {\vee} & \text {by hypothesis} \\ = \mathrm{L} _ {u _ {2} * u _ {1}} ^ {\vee \vee} & \text {(Proposition 23)} \\ = \mathrm{L} _ {(u _ {1} * u _ {2}) \vee} & \text {(Proposition 7)} \end{array}
\]

and hence the proposition is true for \(u_{1} * u_{2}\) . It therefore suffices to prove the proposition when \(u \in \mathrm{T}_{e}(\mathrm{G})\) . Now \(\mathbf{L}_u\) is defined by G operating on G on the right (no. 6) and hence \(\theta_{\mathrm{L}_u}\omega = 0\) since \(\omega\) is right invariant (Differentiable and Analytic Manifolds, R, 8.4.5); therefore, if \(f\) is an analytic function in an open neighbourhood of \(e\) with values in K, then \(\theta_{\mathrm{L}_u}(f\omega) = (\theta_{\mathrm{L}_u}f)\omega\) (Differentiable and Analytic Manifolds, R, 8.4.8). Using the identifications made and Differentiable and Analytic Manifolds, R, 14.4.1, the transpose of \(\mathbf{L}_u\) is \(-\mathbf{L}_u\) , that is \(\mathbf{L}_u^\vee\) .

COROLLARY. Let G be a finite-dimensional real Lie group, \(\mu\) (resp. \(\nu\)) a left (resp. right) Haar measure on G, k an integer \(\geqslant0\) , \(u\in\mathrm{U}_{k}(\mathrm{G})\) and f and g real-valued functions of class \(C^{k}\) on G with compact support. Then

\[
\int_ {G} \left(R _ {u} f\right) g d \mu = \int_ {G} f (R _ {u} ^ {\vee} g) d \mu
\]

\[
\int_ {\mathbf {G}} \left(\mathrm{L} _ {u} f\right) g d \nu = \int_ {\mathbf {G}} f (\mathrm{L} _ {u} ^ {\vee} g) d \nu .
\]

This follows from Proposition 57 and Differentiable and Analytic Manifolds, R, 14.3.8.

